% =====================================================================
% Section 4 — Experimental Results
%
% PROVENANCE. Every quantitative statement in this section resolves to a
% table or figure, and every table or figure resolves to the immutable
% registry below. Nothing here is quoted from an earlier report.
%
%   Registry   : docs/reports/experiments/phase_1b/final_b_new_recomputed/
%                20260819T043733Z/
%   Created    : 2026-08-19T04:37:33.586402+00:00
%   Git commit : 75df6c061eca5a3d63cfac706d1b6e256e404342
%   Split      : task-B-group-split:359be7f5379818d4
%
%   tab:task_b           <- main_test_summary.csv        (15 runs)
%   tab:ablation         <- ablation_vs_main_seed42.csv  (20 runs)
%   tab:perspecies_full  <- per_class_metrics.csv        (15 runs)
%   fig:confusion        <- predictions.csv              (3 runs)
%   fig:perclass         <- per_class_metrics.csv        (15 runs)
%
% NGOAI REGISTRY (nguon khac, lineage rieng -- ghi ro trong tung caption):
%   tab:efficiency  <- results/efficiency.csv
%                      thuoc tinh kien truc, do voi pretrained=False;
%                      latency da bo (do tren RTX 3060, khac phan cung registry)
%   Sec 4.5 probe   <- docs/reports/experiments/phase_1b/anatomy_destruction_probe/
%                      tables/anatomy_probe_summary.csv + run_summary.json
%   Sec 4.5 hoi tu  <- data/experiments/Task-B-new/final/**/test_metrics.json
%                      (best_epoch / epochs_trained / stopped_early / final acc)
%   fig:gradcam     <- artifact/person1/gradcam/  (repo WoodID, branch p1,
%   fig:tsne        <- artifact/person1/tsne/      commit d39a3f6; xem SOURCE.txt)
%                      ca hai dung checkpoint seed 42 cua CUNG registry nay
%
% Per-table run-id lists are carried in artifact/person3/tables/
% tables_metadata.json and in the header comment of each .tex file.
%
% REPORTING RULE. Dispersion is always a standard deviation over the three
% main seeds. It is never an interval estimate, and no claim of statistical
% reliability is attached to any difference in this section.
% =====================================================================

\section{Experimental Results}
\label{sec:results}

\subsection{Experimental Setup}
\label{sec:results:setup}

All runs reported in this section come from a single immutable evaluation
registry, created on 2026-08-19 against git commit
\texttt{75df6c06} and split version
\texttt{task-B-group-split:\allowbreak 359be7f5379818d4}. Re-running the evaluation
notebook against that commit reproduces every number below.

\paragraph{Data.} The Task~B partition contains 17 conifer species split
3{,}365 / 903 / 532 images across train / validation / test. Three further
species present in the source collection
(\textit{Calocedrus} sp., \textit{Cupressus funebris},
\textit{Pinus} sp.) were excluded because they are identified only to genus.
The partition is grouped by \emph{origin key}: images sharing an origin key
are assigned to a single split, rather than being divided image-by-image. The
origin key is derived from the filename and is intended to stand for the
physical wood specimen, so that test images come from specimens the model has
not seen rather than from unseen crops of familiar specimens. Whether the key
corresponds one-to-one to a physical specimen has not been established, and
Section~\ref{sec:discussion:originkey} sets out what the grouping does and
does not guarantee.

\paragraph{Family assignment.} Species are grouped into five families ---
Araucariaceae, Cupressaceae, Pinaceae, Podocarpaceae and Taxaceae --- following
accepted conifer taxonomy, so that \textit{Araucaria araucana} is assigned to
Araucariaceae. This differs from the four-family scheme of the project's
family-level task, which folds that species into Pinaceae. Family membership
affects only how results are grouped for presentation; it plays no part in
training or evaluation, and every per-species number is unchanged by the
choice. Family-level error rates in this section are computed under the
five-family scheme throughout.

\paragraph{Training configuration.} All five backbones were initialised from
ImageNet weights and fine-tuned for 50 epochs with a cosine schedule, a head
learning rate of $1\times10^{-4}$, a backbone learning rate of
$1\times10^{-5}$, weight decay $1\times10^{-4}$, label smoothing 0.1, and
mixed-precision training. These settings are identical across all 30
registered runs. Batch size is the one setting that varies, and it varies by
backbone rather than by experiment: 64 for ResNet-50 and MobileNetV3-Large,
32 for EfficientNet-B3, ConvNeXt-Tiny and Swin-Tiny, held constant for a
given backbone across the main runs and all three ablations. Comparisons
between an ablation and its baseline are therefore within-backbone and
single-variable; comparisons \emph{between} backbones carry batch size as an
additional uncontrolled difference.

\paragraph{Seeds and evaluation.} The main configuration was trained with
three seeds (42, 123, 456), giving 15 main runs. Each ablation was trained
with seed 42 only, giving 5 runs per ablation. Evaluation used FP16
autocast, TF32 matrix multiplication and channels-last memory format on an
NVIDIA GeForce RTX 4090 under PyTorch 2.11.0+cu128, with the whole split
resident on the GPU so that batching does not affect the result.

\subsection{Main Benchmark}
\label{sec:results:benchmark}

Table~\ref{tab:task_b} reports test accuracy and macro-F1 for the five
architectures, as the mean and standard deviation over the three main seeds.

The two modern backbones lead. ConvNeXt-Tiny reaches
$99.561 \pm 0.354$\,\% accuracy and $99.596 \pm 0.326$\,\% macro-F1,
ahead of Swin-Tiny at $99.123 \pm 0.581$\,\% and
$99.196 \pm 0.533$\,\%. MobileNetV3-Large follows at
$98.371 \pm 0.319$\,\% accuracy despite being the smallest model in the
comparison, and EfficientNet-B3 at $97.807 \pm 0.319$\,\%. ResNet-50 is
last by a clear margin, $92.356 \pm 1.022$\,\% accuracy and
$92.569 \pm 0.982$\,\% macro-F1 --- roughly seven points below the leader
and with the largest seed-to-seed spread of the five.

Two features of the spread are worth stating before any ranking is read too
firmly. First, the gaps separating ConvNeXt-Tiny, Swin-Tiny and
MobileNetV3-Large are of the same order as the seed-to-seed standard
deviations themselves, so the ordering among those three rests on
differences that a different choice of seeds could plausibly reorder.
Second, ResNet-50's separation from the rest is several times its own
standard deviation and is the one ordering in the table that does not depend
on that caveat.

Accuracy and macro-F1 track each other closely for every backbone, which is
expected given that the test partition is close to balanced --- per-species
support ranges from 24 to 39 images (Table~\ref{tab:perspecies_full}). No
backbone is buying accuracy by neglecting a small class.

\subsection{Per-Species and Confusion Analysis}
\label{sec:results:perspecies}

\paragraph{Model selection.} The confusion matrix is reported for a single
architecture, selected by a criterion fixed before the confusion structure
was inspected: the backbone with the highest mean macro-F1 over the three
main seeds, which is ConvNeXt-Tiny (Table~\ref{tab:task_b}). The
alternative of reporting whichever individual run scored highest was
rejected, and is in any case ill-defined here: seeds 42 and 123 tie exactly
at 99.812\,\% accuracy with no stated tie-break.

\paragraph{Aggregation.} Figure~\ref{fig:confusion} sums the three per-seed
$17\times17$ count matrices element-wise, giving $532\times3 = 1{,}596$
predictions. Each seed contributes equally, and row normalisation divides by
the row total of the summed counts. The per-seed matrices are published
alongside the aggregate so the sum can be checked.

\paragraph{Result.} ConvNeXt-Tiny makes 7 errors in 1{,}596 predictions, and
all 7 are the same confusion: \textit{Pinus krempfii} predicted as
\textit{Pinus dalatensis}. The matrix is otherwise exactly diagonal. Both
species are Pinaceae, so no error crosses a family boundary.

Figure~\ref{fig:perclass} and Table~\ref{tab:perspecies_full} give per-species
F1 for all five backbones. Three patterns stand out.

\textit{First}, the \textit{P.\ krempfii} / \textit{P.\ dalatensis} pair is
the only confusion that appears in every architecture: 18, 12, 15, 9 and 7
occurrences for ResNet-50, EfficientNet-B3, MobileNetV3-Large, Swin-Tiny and
ConvNeXt-Tiny respectively, 61 in total, against a single occurrence in the
reverse direction. It accounts for 61 of the 204 errors made by all
backbones and seeds combined, and it does not diminish as capacity
increases. Both species are \textit{Pinus} (Pinaceae) with comparable test
support (36 and 32 images), so neither class imbalance nor a family boundary
explains it. These two species also carry the largest seed-to-seed standard
deviation within the two strongest columns --- 3.30 to 4.08 points, against
zero for every other species under ConvNeXt-Tiny --- so the magnitude of the
effect is itself unstable.

\textit{Second}, the errors are not evenly distributed across
architectures. ResNet-50 alone produces 122 of the 204 errors. Its worst
case is \textit{Podocarpus pilgeri}, at $61.994 \pm 9.598$\,\% F1 ---
33.4 points below the next-worst model and the largest dispersion of any
cell in Table~\ref{tab:perspecies_full} --- driven by 48 confusions with
\textit{Chamaecyparis obtusa}. The remaining four architectures score
between 95.420 and 100.000\,\% F1 on the same species, and two of them make
this error zero times. Pooling architectures would present this as the
second-largest difficulty in the dataset; disaggregated, it is one model's
failure mode on this configuration and is not evidence that the two woods
resemble one another.

\textit{Third}, family-level difficulty does not follow the intuition that
the largest family is the hardest. Error rates by true family are
1.56\,\% for Araucariaceae, 0.78\,\% for Cupressaceae, 2.83\,\% for
Pinaceae, 10.93\,\% for Podocarpaceae and 7.69\,\% for Taxaceae.
Cupressaceae, the largest family at eight species, is the \emph{lowest} of
the five, and all eight of its species are error-free under ConvNeXt-Tiny.
Three of the five families contain exactly one species (Araucariaceae,
Podocarpaceae, Taxaceae), so their rates are single-species rates over much
smaller denominators and are not comparable in kind with the eight-species
and six-species rates. Excluding ResNet-50 changes the picture again:
Araucariaceae falls to 0.00\,\%, Podocarpaceae to 2.08\,\% and Taxaceae to
3.53\,\%, while Cupressaceae reaches 0.49\,\% and Pinaceae 2.13\,\%. Three
of the five families therefore owe most of their apparent difficulty to a
single backbone, and the ordering among the remaining families is not stable
between the two groupings.

\subsection{Rotation and Label Smoothing}
\label{sec:results:ablation}

Table~\ref{tab:ablation} reports three ablations against the main
configuration at the same seed and the same backbone: discrete
$90^\circ$ rotation in place of continuous $\pm30^\circ$ rotation, and label
smoothing set to 0.0 and 0.2 in place of 0.1. Each ablation was trained with
seed 42 only, so every entry is a single-run point estimate with no
dispersion of its own.

Discrete $90^\circ$ rotation does not produce a consistent effect. Its
accuracy deltas are $+0.940$, $+1.316$, $0.000$, $-0.564$ and $-0.564$
percentage points for ResNet-50, EfficientNet-B3, MobileNetV3-Large,
ConvNeXt-Tiny and Swin-Tiny --- positive for the two weaker convolutional
backbones, zero for one, and negative for both of the strongest models. The
sign is not stable across architectures, and four of the five magnitudes lie
within about 1.6 times the corresponding seed-to-seed standard deviation
from Table~\ref{tab:task_b}. Only EfficientNet-B3 shows a delta clearly
larger than its own seed noise, at roughly four times that standard
deviation, and with one seed it cannot be separated from a favourable draw.

Removing label smoothing entirely ($0.0$) lowers accuracy for all five
backbones: $-0.376$, $-0.752$, $-0.376$, $-0.188$ and $-0.188$ percentage
points. The individual magnitudes are small --- three of the five are below
the corresponding seed-to-seed standard deviation --- but the direction is
the same in all five cases, which is the only consistent pattern in the
ablation table. Raising label smoothing to $0.2$ produces no such
consistency: two backbones improve or stay level and three decline.

Taken together, neither ablation supports a recommendation. The rotation
result argues against treating discrete $90^\circ$ augmentation as a general
improvement for this task, and the label-smoothing sweep suggests that the
default of 0.1 is at least not worse than either neighbour tested. Both
readings are limited by the single-seed design
(Section~\ref{sec:discussion:oneseed}).

\subsection{Training Convergence and Anatomical Sufficiency}
\label{sec:results:convergence}

\paragraph{Convergence.} Per-epoch curves were not persisted for the registry
runs, so convergence is characterised from the summary fields recorded with
each checkpoint. The two strongest backbones converge in roughly a third of
the budget: the best validation epoch averages 20.3 for ConvNeXt-Tiny and
15.0 for Swin-Tiny, against 40.3--44.0 for MobileNetV3-Large,
EfficientNet-B3 and ResNet-50. Early stopping fired in all three seeds for
ConvNeXt-Tiny and Swin-Tiny and in none for ResNet-50. The gap between final
training and final validation accuracy separates the same two groups: 0.17
and 0.48 points for ConvNeXt-Tiny and Swin-Tiny against 3.14, 3.85 and 4.73
for the other three. The ordering by convergence speed and by generalisation
gap therefore matches the ordering by test accuracy in
Table~\ref{tab:task_b}.

\paragraph{Anatomical sufficiency.} A control experiment tests whether the
task actually requires anatomical detail. Three probes were trained on the
same partition and evaluated once on the same test set: \textbf{B1}, images
downsampled to $16\times16$ and resized back to $224\times224$; \textbf{B2},
images blurred with a Gaussian kernel of size 49 and $\sigma=8$; and
\textbf{B3}, a small multi-layer perceptron given only a normalised
$3\times32$-bin RGB histogram, discarding all spatial structure. B1 and B2
use a freshly initialised ImageNet-pretrained Swin-Tiny; the full-image
reference is the corresponding Swin-Tiny run at 99.060\,\% accuracy.

\begin{center}
\small
\begin{tabular}{llcc}
\toprule
\textbf{Probe} & \textbf{Input} & \textbf{Accuracy (\%)} & \textbf{Retention (\%)} \\
\midrule
Uniform chance      & none                        & 5.882  & --- \\
Majority class      & none                        & 7.331  & --- \\
Full image          & $224\times224$              & 99.060 & 100.0 \\
B1                  & $16\times16$ upsampled      & 94.361 & 95.3 \\
B2                  & Gaussian blur $\sigma=8$    & 95.677 & 96.6 \\
B3                  & RGB histogram only          & 37.970 & 38.3 \\
\bottomrule
\end{tabular}
\end{center}

At $16\times16$ every structure a wood anatomist would use --- tracheid
diameter, resin canals, ray cells, growth-ring boundaries --- is destroyed,
yet accuracy retains 95.3\,\% of the full-image reference. Heavy blurring
retains 96.6\,\%. Colour alone is not sufficient: B3 reaches 37.970\,\%,
well above the 5.882\,\% chance level but far below either spatial probe.
The task as posed is therefore largely solvable from low-frequency image
content that survives the destruction of anatomical detail, and the
$99$\,\%-level accuracies in Table~\ref{tab:task_b} cannot be read as
evidence that the models perform anatomical identification.
Section~\ref{sec:discussion:shortcut} takes up what this implies.

\subsection{Model Complexity}
\label{sec:results:efficiency}

Table~\ref{tab:efficiency} reports parameter count, multiply--accumulate
operations and model size at the $224\times224$ input resolution.

The accuracy ranking does not follow model size. MobileNetV3-Large reaches
$98.371 \pm 0.319$\,\% accuracy (Table~\ref{tab:task_b}) with 4.22\,M
parameters and 0.22\,G MACs --- roughly one sixth the parameters and one
twentieth the compute of ConvNeXt-Tiny, for a difference of 1.19 accuracy
points. ResNet-50, at 23.54\,M parameters and 4.13\,G MACs, is within 15\,\%
of ConvNeXt-Tiny's compute yet scores seven points lower, so the gap between
them is not explained by capacity. ConvNeXt-Tiny and Swin-Tiny are nearly
identical in complexity (27.83 versus 27.53\,M parameters, 4.46 versus
4.37\,G MACs), which makes their accuracy difference a comparison of
architecture rather than of scale.

Two limits apply to this table. Inference latency is not reported: the
available measurements were taken on different hardware from the training
and evaluation runs, so they cannot be placed alongside any other number in
this paper. And the complexity figures are architecture properties measured
on randomly initialised models, so unlike every other table here they do not
derive from the evaluation registry.

\subsection{Grad-CAM and Feature-Space Structure}
\label{sec:results:gradcam}

\paragraph{Grad-CAM.} Figure~\ref{fig:gradcam} shows Grad-CAM maps for
ConvNeXt-Tiny, Swin-Tiny and MobileNetV3-Large on five test images. The
images were selected by a rule fixed in advance: among images classified
correctly by all three models, the one nearest the median ConvNeXt-Tiny
confidence for its family, with ties broken by image path. No image was
chosen for the appearance of its heatmap. All three maps come from the
seed-42 checkpoints of the same runs used in
Table~\ref{tab:task_b}.

The attributions are diffuse and distributed across the field rather than
concentrated on identifiable anatomical landmarks, and the three
architectures highlight different regions of the same image. Read together
with the probes of Section~\ref{sec:results:convergence}, this is what would
be expected if the decisive evidence were low-frequency and spread over the
image rather than carried by localised anatomical features. Grad-CAM shows
where gradient support lies; it does not by itself establish what feature is
being used, and no anatomical reading of these maps is offered here.

\paragraph{Feature space.} Figure~\ref{fig:tsne} embeds the 768-dimensional
pre-logit features of ConvNeXt-Tiny seed 42 for all 532 test images, reduced
by PCA to 50 dimensions and then by t-SNE to two, with perplexity 30 and a
fixed random seed. Species form largely separated clusters, consistent with
the near-diagonal confusion matrix of Figure~\ref{fig:confusion}. The one
structure that matters for interpretation is the relationship between
\textit{Pinus krempfii} and \textit{Pinus dalatensis}, the single pair that
survives in the best model (Section~\ref{sec:results:perspecies}).

The embedding is a projection chosen to preserve local neighbourhoods, and
distances in it are not metric; it is included as a qualitative view of the
representation, not as evidence of separability.
